\documentclass[12pt,a4paper]{article}

\usepackage[margin=2.5cm]{geometry}
\usepackage{amsmath}
\usepackage{amssymb}
\usepackage{booktabs}
\usepackage{array}
\usepackage{float}
\usepackage{setspace}
\usepackage{lmodern}

\onehalfspacing

\title{\textbf{Primary Econometric Analysis}}
\author{ECC3479}
\date{}

\begin{document}

\maketitle

% ─────────────────────────────────────────────────────────────────────────────
\section{Variable Definitions and Supporting Formulas}
% ─────────────────────────────────────────────────────────────────────────────

All dependent and independent variables are expressed as year-on-year percentage
changes computed from the underlying index series. Let $P^{h}_{t}$ denote the
headline CPI index, $P^{e}_{t}$ the energy CPI index, and $C_{t}$ the domestic
consumption expenditure index (all rebased to 2015\,=\,100). The YoY percentage
changes used in estimation are defined as:

\begin{equation}
    \pi^{h}_{t} = \frac{P^{h}_{t} - P^{h}_{t-12}}{P^{h}_{t-12}} \times 100
    \qquad \text{(monthly, Regressions 1, 3, 4)}
\end{equation}

\begin{equation}
    \pi^{e}_{t} = \frac{P^{e}_{t} - P^{e}_{t-12}}{P^{e}_{t-12}} \times 100
    \qquad \text{(monthly; resampled to quarterly mean for Regression 2)}
\end{equation}

\begin{equation}
    \Delta c_{t} = \frac{C_{t} - C_{t-4}}{C_{t-4}} \times 100
    \qquad \text{(quarterly, Regression 2)}
\end{equation}

The post-shock dummy and the EPG policy dummy are defined as indicator functions:

\begin{equation}
    D^{\text{shock}}_{t} =
    \begin{cases}
        1 & \text{if } t \geq \text{October 2022} \\
        0 & \text{otherwise}
    \end{cases}
\end{equation}

\begin{equation}
    D^{\text{EPG}}_{t} =
    \begin{cases}
        1 & \text{if } \text{October 2022} \leq t \leq \text{June 2023} \\
        0 & \text{otherwise}
    \end{cases}
\end{equation}

The cumulative energy pass-through coefficient, reported for each regression, is
the sum of all estimated energy CPI coefficients across contemporaneous and
lagged terms:

\begin{equation}
    \hat{\beta}^{\text{cum}} = \sum_{k=0}^{3} \hat{\beta}_{k}
\end{equation}

All standard errors are computed using the Newey-West HAC variance estimator.
For a regression with $n$ observations, regressor matrix $\mathbf{X}$, and
residuals $\hat{\varepsilon}$, the HAC variance matrix is:

\begin{equation}
    \widehat{\text{Var}}_{\text{HAC}}(\hat{\boldsymbol{\beta}})
    = (\mathbf{X}^{\top}\mathbf{X})^{-1}
    \hat{\boldsymbol{\Omega}}
    (\mathbf{X}^{\top}\mathbf{X})^{-1}
\end{equation}

where the long-run variance matrix $\hat{\boldsymbol{\Omega}}$ is estimated as:

\begin{equation}
    \hat{\boldsymbol{\Omega}} = \hat{\boldsymbol{\Gamma}}_{0}
    + \sum_{j=1}^{m} w_{j}
    \bigl(\hat{\boldsymbol{\Gamma}}_{j} + \hat{\boldsymbol{\Gamma}}_{j}^{\top}\bigr),
    \qquad
    w_{j} = 1 - \frac{j}{m+1},
    \qquad
    \hat{\boldsymbol{\Gamma}}_{j} = \frac{1}{n}\sum_{t=j+1}^{n}
    \mathbf{x}_{t}\mathbf{x}_{t-j}^{\top}\hat{\varepsilon}_{t}\hat{\varepsilon}_{t-j}
\end{equation}

with Bartlett kernel weights $w_{j}$ and maximum lag truncation $m = 3$ in all
regressions.

% ─────────────────────────────────────────────────────────────────────────────
\section{Econometric Specification}
% ─────────────────────────────────────────────────────────────────────────────

\subsection{Regression 1 — Energy CPI Pass-Through to Headline CPI}

\textbf{Functional form.}
A distributed lag (DL) model with an autoregressive term, estimated separately
for each country by OLS:

\begin{equation}
    \pi^{h}_{t} = \alpha + \sum_{k=0}^{3} \beta_{k}\,\pi^{e}_{t-k}
    + \lambda\,\pi^{h}_{t-1} + \gamma\,D^{\text{shock}}_{t} + \varepsilon_{t}
\end{equation}

where $\pi^{h}_{t}$ is the headline CPI year-on-year percentage change at time
$t$ and $\pi^{e}_{t-k}$ is the energy CPI year-on-year percentage change at lag
$k$.

\textbf{Regressors.}
\begin{itemize}
    \item $\pi^{e}_{t},\,\pi^{e}_{t-1},\,\pi^{e}_{t-2},\,\pi^{e}_{t-3}$ —
    contemporaneous and three monthly lags of energy CPI YoY~\%, capturing
    delayed pass-through via transport costs, food prices, and services.
    \item $\pi^{h}_{t-1}$ — one-period lag of the dependent variable (AR(1)
    term), controlling for the high persistence of headline inflation.
    \item $D^{\text{shock}}_{t}$ — binary post-shock dummy equal to one from
    October 2022 onwards, testing whether the shock produced inflation beyond
    what the energy price series alone predicts.
\end{itemize}

The cumulative pass-through over the four-period horizon is:

\begin{equation}
    \hat{\beta}^{\text{cum}} = \hat{\beta}_{0} + \hat{\beta}_{1}
    + \hat{\beta}_{2} + \hat{\beta}_{3}
\end{equation}

\textbf{Sample.}
Monthly observations. Canada and USA: April 1988 -- March 2025 ($n = 444$).
United Kingdom: April 1989 -- March 2025 ($n = 432$). The binding constraint is
the availability of both the energy CPI series and the headline CPI series with
sufficient lags.

\textbf{Error structure.}
HAC (Heteroscedasticity and Autocorrelation Consistent) standard errors using the
Newey-West estimator with a maximum lag truncation of three months. This corrects
for serial correlation in the residuals and for heteroscedasticity arising from
time-varying volatility in energy prices, without imposing a parametric model on
the error process.

\textbf{Results.}
Table~\ref{tab:reg1} reports the estimated coefficients for each country.

\begin{table}[H]
\centering
\caption{Regression 1 — Energy CPI Pass-Through to Headline CPI}
\label{tab:reg1}
\begin{tabular}{lcccccc}
\toprule
 & \multicolumn{2}{c}{\textbf{Canada}} & \multicolumn{2}{c}{\textbf{United Kingdom}} & \multicolumn{2}{c}{\textbf{United States}} \\
\cmidrule(lr){2-3}\cmidrule(lr){4-5}\cmidrule(lr){6-7}
 & Coef. & $p$ & Coef. & $p$ & Coef. & $p$ \\
\midrule
Constant                     & 0.040  & 0.162 & 0.021  & 0.376 & 0.006  & 0.788 \\
$\pi^{e}_{t}$                & 0.089  & 0.000 & 0.066  & 0.000 & 0.088  & 0.000 \\
$\pi^{e}_{t-1}$              & $-$0.080 & 0.000 & $-$0.060 & 0.000 & $-$0.084 & 0.000 \\
$\pi^{e}_{t-2}$              & 0.002  & 0.599 & $-$0.001 & 0.804 & 0.003  & 0.292 \\
$\pi^{e}_{t-3}$              & 0.000  & 0.990 & 0.002  & 0.509 & $-$0.004 & 0.112 \\
$\pi^{h}_{t-1}$              & 0.965  & 0.000 & 0.978  & 0.000 & 0.994  & 0.000 \\
$D^{\text{shock}}_{t}$       & $-$0.026 & 0.604 & 0.004  & 0.951 & $-$0.120 & \textbf{0.003} \\
\midrule
$\hat{\beta}^{\text{cum}}$   & \multicolumn{2}{c}{$+$0.012} & \multicolumn{2}{c}{$+$0.007} & \multicolumn{2}{c}{$+$0.004} \\
$R^{2}$                      & \multicolumn{2}{c}{0.969}    & \multicolumn{2}{c}{0.989}    & \multicolumn{2}{c}{0.991}    \\
$n$                          & \multicolumn{2}{c}{444}      & \multicolumn{2}{c}{432}      & \multicolumn{2}{c}{444}      \\
\bottomrule
\multicolumn{7}{l}{\footnotesize HAC standard errors (Newey-West, $m=3$). Bold: significant at 5\%.}
\end{tabular}
\end{table}

All three models achieve very high fit ($R^{2} \geq 0.97$), driven by the AR(1)
term: inflation is highly persistent and the previous month's rate explains
nearly all variation in the current month's. The energy CPI coefficients display
a consistent pattern across all three countries — a significant positive effect
at lag~0 that is almost entirely reversed at lag~1, with lags~2 and~3
statistically indistinguishable from zero. This indicates that the energy price
signal is incorporated into headline CPI within one month, with no sustained
multi-period propagation through the distributed lag structure. The net
cumulative pass-through is therefore small: Canada $+0.012$, United Kingdom
$+0.007$, United States $+0.004$.

The post-shock dummy is the key discriminating result. For Canada and the United
Kingdom the coefficient is small and entirely insignificant ($p = 0.604$ and
$p = 0.951$ respectively), meaning neither country experienced headline inflation
materially above the level predicted by the energy price series alone after
October 2022. For the United Kingdom this is consistent with the Energy Price
Guarantee compressing consumer-facing energy prices and thereby suppressing the
inflation signal that would otherwise have appeared in $\pi^{e}_{t}$. For Canada
it reflects the partial insulation afforded by domestic energy production. For
the United States, by contrast, the post-shock dummy is negative and strongly
significant ($\hat{\gamma} = -0.120$, $p = 0.003$): headline CPI ran
approximately 0.12 percentage points \textit{lower} per month than the
historical energy--inflation relationship would predict. This is consistent with
the structural decoupling of US domestic natural gas prices (Henry Hub) from
European TTF spot prices following the shale revolution, meaning the 2022 shock
— which was primarily a European supply event — had a more limited impact on US
consumer energy costs than the global energy CPI measure implies.

% ─────────────────────────────────────────────────────────────────────────────

\subsection{Regression 2 — Energy CPI Pass-Through to Domestic Consumption}

\textbf{Functional form.}
The same distributed lag structure as Regression~1, applied to domestic
consumption growth:

\begin{equation}
    \Delta c_{t} = \alpha + \sum_{k=0}^{3} \beta_{k}\,\pi^{e}_{t-k}
    + \lambda\,\Delta c_{t-1} + \gamma\,D^{\text{shock}}_{t} + \varepsilon_{t}
\end{equation}

where $\Delta c_{t}$ is the domestic consumption year-on-year percentage change
(derived from private final consumption expenditure indices rebased to 2015=100).

\textbf{Regressors.}
\begin{itemize}
    \item $\pi^{e}_{t},\,\pi^{e}_{t-1},\,\pi^{e}_{t-2},\,\pi^{e}_{t-3}$ —
    contemporaneous and three quarterly lags of energy CPI. The energy CPI
    series is resampled from monthly to quarterly frequency by taking the
    quarterly mean prior to estimation.
    \item $\Delta c_{t-1}$ — one-period lag of consumption growth, controlling
    for expenditure persistence.
    \item $D^{\text{shock}}_{t}$ — post-shock dummy as defined above.
\end{itemize}

\textbf{Sample.}
Quarterly observations. The sample is shorter than Regression~1 and differs by
country, constrained by the availability of consumption data. Canada: October
1982 -- July 2023 ($n = 164$). United Kingdom: October 1996 -- July 2023
($n = 108$). United States: October 2008 -- January 2025 ($n = 66$). All
consumption YoY~\% changes are computed as $(\text{index}_{t} /
\text{index}_{t-4} - 1) \times 100$ to maintain a consistent quarterly
year-on-year definition.

\textbf{Error structure.}
HAC standard errors (Newey-West, maximum lag three quarters). The choice of
three lags is consistent with the quarterly frequency and the DL order of the
model.

% ─────────────────────────────────────────────────────────────────────────────

\subsection{Regression 3 — Pooled Panel with Interaction Terms}

\textbf{Functional form.}
A pooled OLS panel regression stacking all three country-time series, with the
United States as the baseline category and country-specific interaction terms to
identify differential pass-through:

\begin{equation}
    \pi^{h}_{it} = \alpha + \beta\,\pi^{e}_{it}
    + \delta_{\text{CAN}}\bigl(\text{CAN}_{i} \times \pi^{e}_{it}\bigr)
    + \delta_{\text{UK}}\bigl(\text{UK}_{i} \times \pi^{e}_{it}\bigr)
    + \mu_{\text{CAN}}\,\text{CAN}_{i} + \mu_{\text{UK}}\,\text{UK}_{i}
    + \lambda\,\pi^{h}_{i,t-1} + \gamma\,D^{\text{shock}}_{t} + \varepsilon_{it}
\end{equation}

The total contemporaneous pass-through for each country is recovered as:

\begin{align}
    \text{USA:}    &\quad \hat{\beta}^{\text{USA}}  = \hat{\beta} \\[4pt]
    \text{Canada:} &\quad \hat{\beta}^{\text{CAN}}  = \hat{\beta} + \hat{\delta}_{\text{CAN}} \\[4pt]
    \text{UK:}     &\quad \hat{\beta}^{\text{UK}}   = \hat{\beta} + \hat{\delta}_{\text{UK}}
\end{align}

A $t$-test on $\hat{\delta}_{\text{CAN}}$ or $\hat{\delta}_{\text{UK}}$ directly
tests the null hypothesis of equal pass-through relative to the USA baseline:

\begin{equation}
    H_{0}: \delta_{j} = 0 \quad \text{vs} \quad H_{1}: \delta_{j} \neq 0,
    \qquad j \in \{\text{CAN},\,\text{UK}\}
\end{equation}

\textbf{Regressors.}
\begin{itemize}
    \item $\pi^{e}_{it}$ — contemporaneous energy CPI YoY~\% (no lags in the
    panel specification, to avoid collinearity with the interaction terms).
    \item $\text{CAN}_{i} \times \pi^{e}_{it}$,\; $\text{UK}_{i} \times
    \pi^{e}_{it}$ — interaction terms between country dummies and energy CPI,
    identifying differential pass-through relative to the USA baseline.
    \item $\text{CAN}_{i}$,\; $\text{UK}_{i}$ — country fixed effects,
    absorbing time-invariant cross-country level differences in headline CPI.
    \item $\pi^{h}_{i,t-1}$ — lagged dependent variable.
    \item $D^{\text{shock}}_{t}$ — post-shock dummy.
\end{itemize}

\textbf{Sample.}
Monthly, pooled across all three countries. The common overlapping period is
April 1989 -- March 2025 ($N = 1{,}320$ country-month observations, 440 per
country after alignment on the binding UK start date).

\textbf{Error structure.}
HAC standard errors (Newey-West, three lags) applied to the pooled residuals.
Note that the panel is not balanced by construction --- all three series run over
the same period --- but the HAC correction does not account for cross-sectional
dependence between country residuals. A more rigorous panel treatment would use
Driscoll-Kraay standard errors, though with only three countries this is rarely
implemented in practice.

% ─────────────────────────────────────────────────────────────────────────────

\subsection{Regression 4 — UK Headline CPI with EPG Policy Dummy}

\textbf{Functional form.}
The baseline DL specification from Regression~1 estimated on UK data only, with
the addition of a policy dummy for the Energy Price Guarantee period:

\begin{equation}
    \pi^{h}_{t} = \alpha + \sum_{k=0}^{3} \beta_{k}\,\pi^{e}_{t-k}
    + \lambda\,\pi^{h}_{t-1} + \gamma\,D^{\text{shock}}_{t}
    + \phi\,D^{\text{EPG}}_{t} + \varepsilon_{t}
\end{equation}

\textbf{Regressors.}
All regressors from Regression~1 for the UK, plus:
\begin{itemize}
    \item $D^{\text{EPG}}_{t}$ — binary dummy equal to one during October 2022
    to June 2023 (the active EPG period), zero otherwise. This captures the
    residual suppression of headline CPI attributable to the government subsidy
    beyond what the (already consumer-facing, already suppressed) energy CPI
    series records. The coefficient $\hat{\phi}$ is expected to be negative if
    the EPG reduced inflation below the level predicted by the energy model.
\end{itemize}

The regression is run twice --- without and with $D^{\text{EPG}}_{t}$ --- to
allow the change in the energy pass-through coefficients $\hat{\beta}_{k}$ and
the cumulative pass-through $\hat{\beta}^{\text{cum}}$ to be assessed before
and after controlling for the policy intervention. The implied EPG-counterfactual
inflation --- what headline CPI would have been absent the subsidy --- is:

\begin{equation}
    \tilde{\pi}^{h}_{t} = \hat{\pi}^{h}_{t} - \hat{\phi}\,D^{\text{EPG}}_{t}
\end{equation}

where $\hat{\pi}^{h}_{t}$ is the fitted value from the model including
$D^{\text{EPG}}_{t}$ and $\hat{\phi} < 0$ implies the EPG suppressed inflation
relative to the counterfactual path.

\textbf{Sample.}
Identical to the UK in Regression~1: April 1989 -- March 2025 ($n = 432$
monthly observations).

\textbf{Error structure.}
HAC standard errors (Newey-West, three lags), identical to Regressions~1--3.

\end{document}
